\section{Literatür verisiyle yazılım uygulamaları}
\label{sec:spss-b12}

\textbf{Amaç ve veri.} \texttt{b12.csv} ile kaynak cinsiyet kategorisi
ve yıl sonu notunun en az 10 olması arasındaki ilişki incelenir
\citep{cortez2008data}. Satırlar F ve M, sütunlar 10 puanın altı ve
10 ve üzeridir. Sıfır hipotezi bu iki kategorik değişkenin bağımsız
olmasıdır; alternatif hipotez ilişki bulunmasıdır. Anlamlılık düzeyi
0,05'tir. Önceden seçilen yöntem, süreklilik düzeltmesi uygulanmayan
Pearson ki-kare testidir; B11'deki ortalamalar farkı testi değildir.

\textbf{Ortak veri.} Karşılaştırmada kullanılan dosya
\nolinkurl{companion/spss/csv/b12.csv} olup 395 satır ve 10 sütun
içerir. \texttt{sex} için 1=F ve 2=M; \texttt{pass10} için
0: $\texttt{g3}<10$, 1: $\texttt{g3}\geq10$ tanımlıdır.
Sıfır notlu 38 kayıt eksik sayılmaz; eşik altı kategorisinde tutulur.
B06'nın seçim göstergeleri ve ağırlığı kullanılmaz. Paylaşılan CSV
kopyası proje verisiyle aynıdır; bu kontrol Excel içe aktarmanın
doğrulandığı anlamına gelmez. Kodları \texttt{companion} klasörünü
içeren paket kökünden çalıştırın; veri sözlüğü için bkz.
sayfa~\pageref{sec:spss-guide}.

\subsection{IBM SPSS uygulaması}
\textbf{Başlamadan önce.} Önce aşağıdaki veri açma bloğunu
\texttt{EXECUTE} satırı dahil çalıştırın. B11'den kalan etkin veri
yerine B12 dosyası yeniden açılır. Filtre, ağırlık ve dosya bölme
kapalı olmalıdır. Dosya açılamazsa \texttt{/FILE} yolunu düzeltip
bu bloğu yeniden çalıştırın; eski etkin veriyle devam etmeyin.

\begin{spssadimlar}
\spssadim{\emph{Analyze $\to$ Descriptive Statistics $\to$ Crosstabs} yolunu açın.}
\spssadim{\texttt{sex} değişkenini \emph{Row(s)}, \texttt{pass10} değişkenini \emph{Column(s)} alanına taşıyın.}
\spssadim{\emph{Statistics} içinde \emph{Chi-square} ve \emph{Phi and Cramér's V} seçeneklerini işaretleyip \emph{Continue} seçin.}
\spssadim{\emph{Cells} içinde \emph{Observed}, \emph{Expected}, \emph{Row} yüzdeleri ve \emph{Standardized} artıkları seçin. Düzeltilmiş standartlaştırılmış artıklar bu karşılaştırmanın ölçüsü değildir.}
\spssadim{\emph{Continue $\to$ OK} ile çıktıyı alın. \emph{Case Processing Summary} tablosunda 395 geçerli ve sıfır eksik kayıt bulunduğunu kontrol edin.}
\spssadim{\emph{Pearson Chi-Square} satırındaki \emph{Asymptotic Significance (2-sided)} sütununu ve beklenen frekans dipnotunu okuyun. \emph{Symmetric Measures} tablosundan Cramér $V$ değerini alın.}
\end{spssadimlar}

{\raggedright
\noindent\textbf{Komutların görevi.} \texttt{TABLES=sex BY pass10}
satır ve sütunu, \texttt{CHISQ PHI} test ve ilişki büyüklüklerini
belirler. \texttt{CELLS=COUNT EXPECTED ROW SRESID}, gözlenen ve
beklenen sayımları, satır yüzdelerini ve Pearson hücre artıklarını
ister.\par}

\par\noindent\begin{minipage}{\linewidth}
\textbf{Uygulama kodu:} \texttt{b12}. Veri açma ve tanımlama:

\begin{lstlisting}[style=prismSPSS]
SET DECIMAL=DOT.
GET DATA /TYPE=TXT
 /FILE='companion/spss/csv/b12.csv'
 /ENCODING='UTF8'
 /ARRANGEMENT=DELIMITED /DELCASE=LINE
 /DELIMITERS="," /QUALIFIER='"'
 /FIRSTCASE=2
 /VARIABLES=id F8.0 school F8.0 sex F8.0 age F8.0
 studytime F8.0 absences F8.0 g1 F8.0 g2 F8.0
 g3 F8.0 pass10 F8.0.
FILTER OFF.
WEIGHT OFF.
SPLIT FILE OFF.
VARIABLE LABELS
 sex 'Kaynak verideki cinsiyet kategorisi'
 pass10 'Yil sonu notu en az 10'.
VALUE LABELS
 sex 1 'F' 2 'M'
 /pass10 0 '10 puanin alti' 1 '10 ve uzeri'.
VARIABLE LEVEL sex pass10 (NOMINAL).
EXECUTE.
\end{lstlisting}
\end{minipage}\par

\par\noindent\begin{minipage}{\linewidth}
\textbf{Frekanslar ve ilişki testi.} Veri açma bloğundan sonra çalıştırın.

\begin{lstlisting}[style=prismSPSS]
FREQUENCIES VARIABLES=sex pass10.
CROSSTABS /TABLES=sex BY pass10
 /FORMAT=AVALUE TABLES
 /STATISTICS=CHISQ PHI
 /CELLS=COUNT EXPECTED ROW SRESID.
\end{lstlisting}
\end{minipage}\par

\begin{outputguidebox}
\textbf{Paylaşılan SPSS tablolarıyla doğrulanan değerler.}
\emph{Case Processing Summary} tablosunda 395 geçerli ve sıfır
eksik kayıt vardır. Çapraz tabloda F satırı $(75,133)$, M satırı
$(55,132)$; satır toplamları 208 ve 187'dir. Eşiği karşılayanların
satır yüzdeleri sırasıyla yüzde 63,9 ve 70,6 görünür.

\emph{Pearson Chi-Square} satırında $\chi^2(1)=1{,}970$ ve
$p=0{,}160$; dipnotta en küçük beklenen frekans 61,54 ve beklenen
frekansı 5'in altında olan hücre sayısı sıfırdır. \emph{Symmetric
Measures} tablosunda Phi ve Cramér $V$ üç basamakla 0,071'dir.

\emph{Continuity Correction} satırındaki $p=0{,}195$ ve Fisher
testinin iki yönlü $p=0{,}165$ değeri farklı yöntemlere aittir;
Pearson sonucunun yazılımlar arasında uyuşmadığı anlamına gelmez.
Bu karşılaştırmada yöntem, sonuçlara göre değiştirilmez.
\end{outputguidebox}

\par\noindent\begin{minipage}{\linewidth}
\subsection{R uygulaması}
\label{sec:r-gercek-b12}

\texttt{correct=FALSE} süreklilik düzeltmesini kapatır. Aşağıdaki
blokları sırayla ve çok satırlı komutları bütünüyle çalıştırın.

\begin{lstlisting}[style=prismR]
veri <- read.csv("companion/spss/csv/b12.csv")
stopifnot(nrow(veri) == 395, ncol(veri) == 10,
  !anyNA(veri), !anyDuplicated(veri$id),
  all(veri$sex %in% c(1, 2)),
  all(veri$pass10 %in% c(0, 1)),
  all(veri$pass10 == as.integer(veri$g3 >= 10)))
tablo <- table(
  sex=factor(veri$sex, levels=c(1, 2),
             labels=c("F", "M")),
  pass10=factor(veri$pass10, levels=c(0, 1),
                labels=c("10_alti", "10_ve_uzeri")))
sonuc <- chisq.test(tablo, correct=FALSE)
print(c(satir=nrow(veri), sutun=ncol(veri),
  eksik=sum(is.na(veri)), sifir_not=sum(veri$g3 == 0)))
print(addmargins(tablo))
\end{lstlisting}
\end{minipage}\par

\textbf{Dosya yolu sorunu yaşarsanız.} İlk satırın yerine aşağıdaki
komutu çalıştırıp \texttt{b12.csv} dosyasını seçin; sonra kontrol
bloğundan devam edin.
\begin{lstlisting}[style=prismR]
veri <- read.csv(file.choose())
\end{lstlisting}
Ham CSV dosyası analiz çıktısı değildir. Özellikle
\texttt{), digits=15)} gibi kapanış satırlarını tek başına çalıştırmayın;
ilgili \texttt{print} komutunun tamamını seçin.

\par\noindent\begin{minipage}{\linewidth}
\textbf{Karşılaştırılacak sayısal çıktılar.} Satır yüzdelerinin
paydaları F için 208, M için 187'dir. \texttt{residuals}, Pearson
hücre artıklarıdır; R'deki \texttt{stdres} ile aynı değildir.

\begin{lstlisting}[style=prismR]
print(100 * prop.table(tablo, margin=1), digits=15)
print(sonuc$expected, digits=15)
print(sonuc$residuals, digits=15)
print(c(
  n=sum(tablo),
  ki_kare=unname(sonuc$statistic),
  df=unname(sonuc$parameter),
  p=sonuc$p.value,
  min_beklenen=min(sonuc$expected),
  bes_alti_hucre=sum(sonuc$expected < 5),
  cramer_v=sqrt(unname(sonuc$statistic)/sum(tablo))
), digits=15)
\end{lstlisting}

\textbf{Çıktıyı okuma.} Paylaşılan R metninde dört satır yüzdesi,
dört beklenen sayım, dört Pearson artığı ve son özetin yedi değeri
yer alır. $\chi^2=1{,}969808570687404$, $p=0{,}160468182571215$
ve $V=0{,}070617682919937$ bulunmuştur. R burada yeniden
çalıştırılmamış; kullanıcının paylaştığı çıktı karşılaştırılmıştır.
\end{minipage}\par

\par\noindent\begin{minipage}{\linewidth}
\subsection{Python uygulaması}
\label{sec:python-b12}

\texttt{correction=False} aynı Pearson testini seçer. Satır ve
sütun sıraları açıkça sabitlenir. Blokları sırayla çalıştırın.

\begin{lstlisting}[style=prismPythonApp]
import pandas as pd
import numpy as np
from scipy import stats

veri = pd.read_csv("companion/spss/csv/b12.csv")
assert veri.shape == (395, 10)
assert not veri.isna().any().any()
assert veri.id.is_unique
assert veri.sex.isin([1, 2]).all()
assert veri.pass10.isin([0, 1]).all()
assert veri.pass10.eq((veri.g3 >= 10).astype(int)).all()
tablo = pd.crosstab(veri.sex, veri.pass10).reindex(
    index=[1, 2], columns=[0, 1], fill_value=0)
tablo.index = ["F", "M"]
tablo.columns = ["10_alti", "10_ve_uzeri"]
ki_kare, p_degeri, serbestlik, beklenen = (
    stats.chi2_contingency(tablo, correction=False))
yuzdeler = 100 * tablo.div(tablo.sum(axis=1), axis=0)
artiklar = (tablo.to_numpy()-beklenen)/np.sqrt(beklenen)
print("satir:", len(veri), "sutun:", veri.shape[1])
print("eksik:", int(veri.isna().sum().sum()))
print("sifir_not:", int((veri.g3 == 0).sum()))
print(tablo)
print("Satir toplamlari:", tablo.sum(axis=1).to_dict())
print("Sutun toplamlari:", tablo.sum(axis=0).to_dict())
\end{lstlisting}
\end{minipage}\par

\par\noindent\begin{minipage}{\linewidth}
\textbf{Ayrıntılı çıktı.} Aşağıdaki yazdırma düzeni, paylaşılan
Python terminal çıktısındaki 13 ondalık basamaklı değerleri verir.

\begin{lstlisting}[style=prismPythonApp]
for baslik, degerler in [
    ("Satir yuzdeleri", yuzdeler.to_numpy()),
    ("Beklenen sayimlar", beklenen),
    ("Pearson hucre artiklari", artiklar)
]:
    print(baslik)
    cerceve = pd.DataFrame(degerler,
        index=tablo.index, columns=tablo.columns)
    print(cerceve.to_string(
        float_format=lambda sayi: f"{sayi:.13f}"))

ozet = dict(n=int(tablo.to_numpy().sum()),
    ki_kare=ki_kare, df=serbestlik, p=p_degeri,
    min_beklenen=beklenen.min(),
    bes_alti_hucre=int((beklenen < 5).sum()),
    cramer_v=np.sqrt(ki_kare/len(veri)))
for ad, deger in ozet.items():
    print(f"{ad}: {deger:.13f}")
\end{lstlisting}

\textbf{Çıktıyı okuma.} Paylaşılan Python ekranında 395 satır,
10 sütun, sıfır eksik değer ve 38 sıfır not vardır. Sayımlar ve
toplamlar SPSS ile aynıdır. $p=0{,}1604681825712$ ve
$V=0{,}0706176829199$ değerlerini ilgili etiketlerle okuyun.
Kod ham kayıtları değiştirmez; tablo ayrı bir nesnedir.
\end{minipage}\par

\Needspace{12\baselineskip}
\subsection{Sonuçların karşılaştırılması ve ortak rapor}
13 Eylül 2026'da paylaşılan SPSS tabloları, R metni ve Python
terminal ekranı karşılaştırılmıştır. R metnindeki 19 sayısal
özet, Python sonuçlarıyla $10^{-12}$ mutlak tolerans içinde
uyumludur. Python kodu proje CSV'siyle ayrıca çalıştırılmıştır;
ekrandaki 13 ondalık basamaklı sonuçlarla eşleşir. SPSS tabloları
gösterilen yuvarlama hassasiyetlerinde aynı sonuçları verir.
Aşağıdaki tablolar çıktıların kitap düzenine uyarlanmış özetidir;
ekran görüntüsü değildir. Ara hesaplar yuvarlanmamış değerlerle yapılır.

\begin{table}[H]
\centering
\caption{B12'de kaynak kategorisine göre 10 puan eşiği}
\label{tab:b12-dogrulanmis-sayim}
\begin{tabular}{@{}lrrrr@{}}
\toprule
Grup & $<10$ & $\geq10$ & Toplam & $\geq10$ yüzdesi \\
\midrule
F & 75 & 133 & 208 & 63,942308 \\
M & 55 & 132 & 187 & 70,588235 \\
\midrule
Toplam & 130 & 265 & 395 & 67,088608 \\
\bottomrule
\end{tabular}
\par\smallskip
{\footnotesize Son sütun ilgili satırın yüzdesidir: F için
$100(133/208)$, M için $100(132/187)$, toplam için $100(265/395)$.
Tam 10 puan, eşiği karşılayan kategoriye dahildir.\par}
\end{table}

\begin{table}[H]
\centering
\caption{B12'de beklenen sayımlar ve Pearson hücre artıkları}
\label{tab:b12-dogrulanmis-artik}
\begin{tabular}{@{}lrr@{}}
\toprule
Hücre & Beklenen sayım ($E$) & Pearson artığı \\
\midrule
F, $<10$ & 68,455696 & $0{,}790968$ \\
F, $\geq10$ & 139,544304 & $-0{,}553997$ \\
M, $<10$ & 61,544304 & $-0{,}834199$ \\
M, $\geq10$ & 125,455696 & $0{,}584276$ \\
\bottomrule
\end{tabular}
\par\smallskip
{\footnotesize $E=(\text{satır toplamı})(\text{sütun toplamı})/395$;
Pearson artığı $(O-E)/\sqrt{E}$'dir. SPSS \texttt{SRESID}, R
\texttt{residuals} ve Python \texttt{artiklar} burada aynı ölçüdür.
Yuvarlanmamış dört artığın kareleri toplamı Pearson ki-kare değerini verir.\par}
\end{table}

\begin{table}[H]
\centering
\caption{B12'de Pearson testi ve ilişki büyüklüğü}
\label{tab:b12-dogrulanmis-test}
\begin{tabular}{@{}lr@{}}
\toprule
Ölçü & Değer \\
\midrule
Geçerli kayıt sayısı & 395 \\
Pearson $\chi^2$ & 1,969809 \\
Serbestlik derecesi & 1 \\
$p$ değeri & 0,160468 \\
En küçük beklenen frekans & 61,544304 \\
Beklenen frekansı 5'in altındaki hücre & 0 \\
Cramér $V$ & 0,070618 \\
\bottomrule
\end{tabular}
\par\smallskip
{\footnotesize Süreklilik düzeltmesi uygulanmamıştır. Bu $2\times2$
tabloda Cramér $V=\sqrt{\chi^2/395}$'tir; yön belirtmez.\par}
\end{table}

\Needspace{12\baselineskip}
\begin{outputguidebox}
Tablo~\ref{tab:b12-dogrulanmis-sayim}, eşiği karşılayanların
örneklemdeki oranlarını betimler. Yüzdeler farklı olsa da
Tablo~\ref{tab:b12-dogrulanmis-test} içindeki
$p=0{,}160468>0{,}05$ nedeniyle bağımsızlık sıfır hipotezi
reddedilmez. Bu, bağımsızlığın kanıtı veya oranların kesin eşitliği
değildir. Cramér $V=0{,}070618$, ilişki büyüklüğünün örneklem
özetidir; nedensel etki ya da açıklanan varyans yüzdesi değildir.

Tablo~\ref{tab:b12-dogrulanmis-artik} içindeki pozitif artık,
bağımsızlık altında beklenenden fazla; negatif artık daha az
gözlem olduğunu gösterir. Bunlar düzeltilmiş artıklar değildir;
her hücreyi ayrı anlamlılık testi olarak raporlamayın.

Beklenen sayımların yeterli olması, öğrencilerin bağımsızlığını
veya iki okulla sınırlı verinin temsil gücünü doğrulamaz. Test okul,
önceki başarı ve diğer karıştırıcılar için düzeltme içermez.
B11'de yıl sonu \emph{ortalamaları}, burada ise 10 puan eşiğini
karşılama \emph{oranları} karşılaştırılır. Farklı sorulara farklı
kararlar çıkması çelişki değildir; notları iki kategoriye ayırmak
puan farklılıklarına ilişkin bilgiyi kaybettirir.
\end{outputguidebox}

\Needspace{10\baselineskip}
\textbf{Raporlama örneği (doğrulanmış çıktılarla).}
``395 kayıtta yıl sonu notu en az 10 olanların oranı F grubunda
$133/208$ (yüzde 63,942308), M grubunda $132/187$ (yüzde 70,588235)
bulunmuştur. Süreklilik düzeltmesiz Pearson testinde kaynak cinsiyet
kategorisi ile eşiği karşılama durumu arasında yüzde 5 düzeyinde
yeterli ilişki kanıtı elde edilmemiştir
($\chi^2(1)=1{,}969809$, $p=0{,}160468$, Cramér $V=0{,}070618$).
En küçük beklenen frekans 61,544304'tür; 5'in altında beklenen
frekans yoktur. Sıfır notlu 38 kayıt analizde tutulmuştur.
Sonuç bağımsızlığı kanıtlamaz. Gözlemlerin bağımsızlığı varsayılmış;
okul ve diğer karıştırıcılar için düzeltme yapılmamıştır. İki okulla
sınırlı veriden nedensel veya geniş bir evrene yönelik sonuç çıkarılmamıştır.''

\Needspace{8\baselineskip}
\subsection{Karşılaştırmalı görev: aynı öğrenciler, farklı hedefler}
\label{sec:b12-karsilastirmali-gorev}

\textbf{Hedef ve teslim.} B11'in yalnız F--M Welch analiziyle B12'nin
Pearson analizini karşılaştırın; B11'deki eşleştirilmiş dönem farkını bu
iki soruya karıştırmayın. Mevcut çıktıları kullanarak aşağıdaki beş
maddeyi yanıtlayın. Yeni eşik veya yeni anlamlılık testi seçmeyin.
Kontrol yanıtını okumadan önce kendi gerekçeli raporunuzu tamamlayın.

\begin{enumerate}
  \item \textbf{Ne değişti?} İki analizde kayıtlar, gruplar, yanıt
  değişkeni, hedeflenen büyüklük ve sıfır hipotezini karşılaştırın.
  G3'ü iki kategoriye ayırmak hangi bilgiyi kaybettirir?
  \item \textbf{Payda ve yön.} F ve M için eşiği karşılama yüzdelerini
  payları ve paydalarıyla gösterin. F eksi M oran farkını yüzde puan
  olarak hesaplayın. Bu fark B11'in not puanı farkıyla aynı şey midir?
  Cramér $V$ bu farkın işaretini verir mi?
  \item \textbf{Hesap kontrolü.} Yuvarlanmamış Pearson hücre
  artıklarının karelerini toplayıp $\chi^2$ ile karşılaştırın.
  En küçük beklenen sayımı ve 5'in altındaki hücre sayısını belirtin.
  Bu kontrol öğrencilerin bağımsızlığını da doğrular mı?
  \item \textbf{Karşılaştırmalı rapor.} İki soruyu, test kararlarını
  ve ilişki/fark özetlerini içeren kısa bir paragraf yazın.
  ``Welch anlamlı, ki-kare anlamsız; yazılımlar çelişiyor ve ikinci
  test grupların eşitliğini kanıtlıyor'' yorumunu düzeltin.
  \item \textbf{Analiz disiplini.} ``Eşiği değiştirip anlamlı sonuç
  bulalım; sıfır notları da çıkaralım'' önerisini değerlendirin.
  Aynı öğrencilerle yapılan iki analizin neden bağımsız iki tekrar
  çalışma olmadığını ve nedensel sonuç vermediğini açıklayın.
\end{enumerate}

\Needspace{8\baselineskip}
\textbf{Gerekçeli kontrol yanıtı.}
\begin{enumerate}
  \item Her iki analiz de aynı 395 öğrenciyi, 208 F ve 187 M kaydını
  ve sıfır notları kullanır. Welch sayısal G3 için
  $H_0:\mu_{G3,F}=\mu_{G3,M}$ hipotezini sınar. Pearson analizinde
  yanıt $\mathbb{1}(G3\geq10)$'dur; sıfır hipotezi kaynak kategorisi
  ile eşik durumunun bağımsızlığıdır. Bu iki gruplu tabloda bu,
  grupların eşiği karşılama olasılıklarının eşitliğine karşılık gelir.
  Aynı eşik kategorisindeki farklı notlar artık ayırt edilmez.
  Değişen kayıtlar değil, soru, yanıtın gösterimi ve yöntemdir.
  \item Yüzdeler $100(133/208)=63{,}942308$ ve
  $100(132/187)=70{,}588235$'tir. F eksi M farkı yaklaşık
  $-6{,}645928$ yüzde puandır; iki payda da 395 değildir.
  Bu betimsel oran farkı, B11'in $-0{,}948092$ not puanı farkından
  farklı bir büyüklüktür. Cramér $V=0{,}070618$ yön belirtmez.
  Burada oran farkı için ayrıca güven aralığı hesaplanmamıştır.
  \item Dört artığın kareleri toplamı
  $\sum (O-E)^2/E=1{,}969808570687404$ olup Pearson istatistiğidir.
  En küçük beklenen sayım yaklaşık 61,544304, 5'in altındaki hücre
  sayısı sıfırdır. Hücre hesaplarının uygunluğu öğrencilerin
  bağımsızlığını veya örneklemin temsil gücünü kanıtlamaz.
  Yuvarlanmış artıklarla oluşan küçük fark, yöntem çelişkisi değildir.
  \item ``Yıl sonu ortalamaları için F eksi M farkı
  $-0{,}948092$ puandır (\%95 GA
  $[-1{,}850732;\ -0{,}045452]$); Welch testinde
  $p=0{,}039577$ ile eşit ortalamalar hipotezi reddedilmiştir.
  Ayrı eşik sorusunda oranlar yüzde 63,942308 ve 70,588235'tir;
  Pearson testinde $\chi^2(1)=1{,}969809$, $p=0{,}160468$ ve
  $V=0{,}070618$ bulunmuş, bağımsızlık hipotezi reddedilmemiştir.
  Kararlar her test için ayrı 0,05 düzeyindedir; sorular farklı
  olduğundan çelişki yoktur. İkinci sonuç eşitliğin veya bağımsızlığın
  kanıtı değildir.'' Bu karar ayrılığı tek başına yöntemlerden
  birinin üstünlüğünü de göstermez.
  \item Sonuca bakarak eşik seçmek önceden tanımlanan soruyu
  değiştirir; yalnız anlamlı sonucu sunmak uygun değildir.
  Sıfırları dışlamak kayıt kapsamını ve hedefi değiştirir; sıfır
  notlar kendiliğinden eksik değer sayılmaz. Aynı öğrencilerin
  yeniden analizi yeni bağımsız kanıt üretmez. Okul, önceki başarı
  ve diğer karıştırıcılar denetlenmediğinden nedensel açıklama
  yapılamaz; iki okulla sınırlı kaynağın temsil gücü ayrıca tartışılır.
\end{enumerate}

\Needspace{8\baselineskip}
\textbf{Değerlendirme ölçütleri (10 puan).} Her görev maddesindeki
iki bileşenin her biri 1 puandır; doğru ve gerekçeli bileşen 1, eksik
veya yanlış bileşen 0 puan alır. Eşdeğer doğru ifadeler ve belirtilen
basamaklara uygun yuvarlama kabul edilir.
\begin{enumerate}
  \item Ortak kayıt/grupları ve değişen yanıtı ayırma (1);
  iki sıfır hipotezini ve kategorileştirmenin bilgi kaybını açıklama (1).
  \item Pay, payda, yüzdeler ve yüzde puan farkını doğru verme (1);
  not puanı farkından ayrımı ve $V$'nin yönsüzlüğünü açıklama (1).
  \item Artık kareleri toplamını ve beklenen sayım kontrolünü
  doğru verme (1); bunları bağımsızlık kanıtından ayırma (1).
  \item İki sorunun sayısal özetleri ve kararlarını doğru
  raporlama (1); çelişki ve eşitlik kanıtı yorumlarını düzeltme (1).
  \item Eşik değiştirme ve sıfırları silme önerilerinin soru/kapsam
  üzerindeki etkisini açıklama (1); bağımsız tekrar ve nedensellik
  iddialarını sınırlılıklarla birlikte reddetme (1).
\end{enumerate}
Yalnız $p$ değerlerini doğru kopyalamak yeterli değildir. Bağımsızlığı
kanıtlanmış sayan veya nedensel iddiayı koruyan yanıt tam puan alamaz.
Bu ölçütler öğretim amaçlıdır; öğrenme etkileri henüz öğrenci
uygulamasıyla değerlendirilmemiştir.